\subsection{Affine Springer representations}
In this subsection we introduce an analog of Springer's $W$-action on $\cohog{*}{\sB_{e}}$ in the affine situation.  

\subsubsection{The affine Weyl group} We view $W$ as a group of automorphisms of the cocharacter lattice $\xcoch(T)$, where $T$ is a fixed maximal torus of $G$. The {\em extended affine Weyl group} $\tilW$ is the semidirect product
\begin{equation*}
\tilW=\xcoch(T)\rtimes W.
\end{equation*}
When $G$ is simply-connected, so that $\xcoch(T)$ is spanned by coroots, $\tilW$ is a Coxeter group with simple reflections $\{s_{0},s_{1}, \cdots,s_{r}\}$ in bijection with the nodes of the extended Dynkin diagram of $G$.  In general, $\tilW$ is a semidirect product of the {\em affine Weyl group} $\Wa=(\ZZ R^{\vee})\rtimes W$ (where $\ZZ R^{\vee}$ is the coroot lattice), which is a Coxeter group, and an abelian group $\Om\cong \xcoch(T)/\ZZ R^{\vee}$. The group $\tilW$ naturally acts on the affine space $\xcoch(T)_{\RR}$ by affine transformations, where $\xcoch(T)$ acts by translations.

\begin{theorem}[Lusztig {\cite{L96}}, Sage {\cite{Sage}}]\label{th:aff action} There is a canonical action of $\tilW$  on $\homog{*}{\sY_{\gamma}}$. 
\end{theorem}

Since $\sY_{\gamma}$ is not of finite type, the $\ell$-adic homology $\homog{*}{\sY_{\gamma}}$ is understood as the inductive limit $\varinjlim_{n}\homog{*}{\sY_{\gamma,n}}$, whenever we present $\sY_{\gamma}$ as a union of projective subschemes $\sY_{\gamma,n}$.

\subsubsection{Sketch of the construction of the $\tilW$-action} We consider only the case $G$ is simply-connected so that $\tilW=\Wa$ is generated by affine simple reflections $s_{0}, \cdots, s_{r}$. For each parahoric  subgroup $\bP\subset LG$ we have a corresponding affine Springer fiber $\tX_{\bP,\gamma}$. For $\bP$ containing a fixed Iwahori subgroup $\bI$, we have a projection $\pi_{\bP,\gamma}:\tY_{\gamma}\to \tX_{\bP,\gamma}$. 

Let $L_{\bP}$ be  the Levi quotient of $\bP$ and $\frl_{\bP}=\Lie L_{\bP}$. We have an evaluation map $\ev_{\bP,\gamma}: \tX_{\bP,\gamma}\to [\frl_{\bP}/L_{\bP}]$ defined as follows. For $[g]\in\Fl_{\bP}$ such that $\Ad(g^{-1})\gamma\in\Lie \bP$, we send the coset $[g]=g\bP$ to the image of $\Ad(g^{-1})\gamma$ under the projection $\Lie \bP\to\frl_{\bP}$. This is well-defined up to the adjoint action of $L_{\bP}$. We have a Cartesian diagram
\begin{equation*}
\xymatrix{\tY_{\gamma}\ar[d]^{\pi_{\bP,\gamma}}\ar[r]^{\ev_{\bI,\gamma}} & [\wt{\frl}_{\bP}/L_{\bP}]\ar[d]^{\pi_{\frl_{\bP}}}\\
\tX_{\bP,\gamma}\ar[r]^{\ev_{\bP,\gamma}} & [\frl_{\bP}/L_{\bP}]}
\end{equation*}
where $\pi_{\frl_{\bP}}$ is the Grothendieck alteration for the reductive group $\frl_{\bP}$. By the Springer theory for $L_{\bP}$, we have a $W(L_{\bP})$-action on the direct image complex $\bR\pi_{\frl_{\bP,*}}\DD$ (where $\DD$ stands for the dualizing complex for $[\wt{\frl}_{\bP}/L_{\bP}]$). By proper base change, we get an action of $W(L_{\bP})$ on $\bR\pi_{\bP,\gamma,*}\DD_{\tY_{\gamma}}$, and hence on $\homog{*}{\sY_{\gamma}}$. Taking a standard parahoric $\bP$ corresponding to the $i$-th node in the extended Dynkin diagram, then $W(L_{\bP})=\jiao{s_{i}}$, and we get an involution $s_{i}$ acting on $\homog{*}{\sY_{\gamma}}$. To check the braid relation between  $s_{i}$ and $s_{j}$ for neighboring nodes $i$ and $j$, we may choose a standard parahoric $\bP$ such that $W(L_{\bP})=\jiao{s_{i},s_{j}}$ and the braid relation holds because $W(L_{\bP})$ acts on $\homog{*}{\sY_{\gamma}}$. This shows that $\tilW$ acts on $\homog{*}{\sY_{\gamma}}$. \qed


Despite the simplicity of the construction of the $\tilW$-action on the homology of affine Springer fibers, the calculation of these actions are quite difficult. One new feature here is that the action of $\tilW$ on $\homog{*}{\sY_{\gamma}}$ {\em may not be semisimple}, as we shall see in the next example.


\subsubsection{An example in $\SL_{2}$}\label{sss:tilW SL2 ram} Consider the case $G=\SL_{2}$ and the element $\gamma=\mat{0}{t^{2}}{t}{0}$. This is a subregular case. The affine Springer fiber $\sY_{\gamma}$ has two irreducible components $C_{0}$ and $C_{1}$ both isomorphic to $\PP^{1}$. Here $C_{0}$ parametrizes chains of lattices $\{\L_{i}\}$ where $\L_{0}$ is the standard lattice $\calO^{2}_{F}$ and $\L_{-1}\subset \L_{0}$ is varying. The other component $C_{1}$ parametrizes chains of lattices $\{\L_{i}\}$ where $\L_{-1}=t\calO_{F}\oplus\calO_{F}$ and $\L_{0}$ is varying. The fundamental classes $[C_{0}], [C_{1}]$ give a natural basis for the top homology group $\homog{2}{\sY_{\gamma}}$. One can show that the action of the affine Weyl group $\tilW=\jiao{s_{0},s_{1}}$ on $\homog{2}{\sY_{\gamma}}$ takes the following form under the basis $[C_{0}]$ and $[C_{1}]$:
\begin{equation*}
s_{0}=\mat{-1}{2}{0}{1}; \quad s_{1}=\mat{1}{0}{2}{-1}.
\end{equation*} 
We see that $\homog{2}{\sY_{\gamma}}$ is a {\em nontrivial extension} of the sign representation of $\tilW$ by the trivial representation spanned by $[C_{0}]+[C_{1}]$. One can also canonically identify the $\tilW$-module $\homog{2}{\sY_{\gamma},\ZZ}$ with the affine coroot lattice of the loop group $LG$.

\subsection{Comments and generalizations}

\subsubsection{Relation with orbital integrals} As we will see in \S\ref{s:orb}, the cohomology and point-counting of affine Springer fibers are closely related to orbital integrals on $p$-adic groups $G(\FF_{q}\lr{t})$. 

\subsubsection{Extended symmetry} The $\tilW$-action on $\homog{*}{\sY_{\gamma}}$ can be extended to an action of the wreath product $\Sym(\xch(T))\rtimes \tilW$. For homogeneous affine Springer fibers $\sY_{\gamma}$ (those admitting a torus action coming from loop rotation), the equivariant cohomology group $\upH^{*}_{\Gm}(\sY_{\gamma})$ admits an action of the {\em graded double affine Hecke algebra},  which is a deformation of $\Sym(\xch(T))\rtimes \tilW$. For details we refer to \cite{OY}. Vasserot and Varagnolo \cite{Vass} \cite{VV} constructed an action of the {\em double affine Hecke algebra} on the $K$-groups of affine Springer fibers.


\subsubsection{The group version} Taking $\gamma\in G(F)$ instead of in $\frg(F)$, one can similarly define the group version of  affine Springer fibers, which we still denote by $\tX_{\gamma}$ with reduced structure $\sX_{\gamma}$. For a $k$-algebras $R$, we have
\begin{equation}\label{gp AFS}
\tX_{\gamma}(R)=\{[g]\in \Gr_{G}(R)|g^{-1}\gamma g\in L^{+}G(R)\}.
\end{equation}
However, in the group version, the definition above admits an interesting generalization. Recall the $L^{+}G$-double cosets in $LG$ are indexed by dominant cocharacters $\l\in\xcoch(T)^{+}$. For $\l\in\xcoch(T)^{+}$ we denote the corresponding double coset by $(LG)_{\l}$, which is the preimage of the Schubert stratum $\Gr_{G,\l}$ under the projection $LG\to \Gr_{G}$. Similarly we may define $(LG)_{\leq \l}$ to be the preimage of the closure $\Gr_{G,\leq \l}$ of $\Gr_{G,\l}$. One can replace the condition $g^{-1}\gamma g\in L^{+}G(R)$ in \eqref{gp AFS} by $g^{-1}\gamma g\in (LG)_{\l}(R)$ or $g^{-1}\gamma g\in (LG)_{\leq\l}(R)$, and take reduced structures to obtain reduced generalized affine Springer fibers $\sX_{\l,\gamma}$ and $\sX_{\leq\l, \gamma}$. We have an open embedding $\sX_{\l,\gamma}\incl\sX_{\leq\l,\gamma}$, whose complement is the union of $\sX_{\mu,\gamma}$ for dominant coweights $\mu\leq\l$. The motivation for introducing $\sX_{\leq\l,\gamma}$ is to give geometric interpretation of orbital integrals of spherical Hecke functions on $G(F)$. 


A.Bouthier has established the fundamental geometric properties of $\sX_{\l,\gamma}$, parallel to Theorem \ref{th:ft quot} and Theorem \ref{th:aff dim}. 

\begin{theorem}[Bouthier \cite{Bouthier}]  Let $\gamma\in G(F)$ be regular semisimple, and let $\l\in\xcoch(T)^{+}$.
\begin{enumerate}
\item The generalized affine Springer fiber $\sX_{\l,\gamma}$ is non-empty if and only if $[\nu_{\gamma}]\leq\l$, where $[\nu_{\gamma}]\in\xcoch(T)^{+}_{\QQ}$ is the Newton point of $\gamma$, see \cite[\S2]{KV}.
\item The ind-scheme $\sX_{\l,\gamma}$ is locally of finite type.
\item We have
\begin{equation*}
\dim\sX_{\l,\gamma}=\jiao{\rho,\l}+\frac{1}{2}(\val_{F}\Delta(\gamma)-c(\gamma))
\end{equation*}
where $\rho\in\xch(T)_{\QQ}$ is half the sum of positive roots, and $\Delta(\gamma)$ and $c(\gamma)$ are defined similarly as in the Lie algebra situation.
\end{enumerate}
\end{theorem}
The proof of this theorem uses the theory of Vinberg semigroups, which is a kind of compactification of $G$.


\subsubsection{} In \cite{KV}, Kottwitz and Viehmann defined two generalizations of affine Springer fibers for elements  $\gamma$ in the Lie algebra $\frg(F)$.

\subsubsection{} As an analog of Hessenberg varieties, one can also consider the following situation. Let $(\rho,V)$ be a linear representation of a reductive group $G$ over $k$. Let $\L\subset V\otimes F$ be an $\calO_{F}$-lattice stable under $G(\calO_{F})$. For $v\in V\otimes F$ we may define a sub-ind-scheme $\tX_{v}$ of $\Gr_{G}$ 
\begin{equation*}
\tX_{\L, v}(R)=\{[g]\in\Gr_{G}(R)|\rho(g^{-1})v\in\L\widehat{\otimes}_{k}R\}.
\end{equation*}
Let $\sX_{\L,v}$ be the reduced structure of $\tX_{\L,v}$. The cohomology of these ind-schemes are related to orbital integrals that appear in relative trace formulae.



\subsection{Exercises}

\subsubsection{}\label{exII:SO parahoric} Let $G=\SO(V,q)$ for some vector space $V$ over $k$ equipped with a quadratic form $q$. Give an interpretation of the parahoric subgroups and affine partial flag varieties of $LG$ in terms of self-dual lattice chains in $V\otimes_{k}F$, in the same style as in \S\ref{sss:Sp parahoric}.

\subsubsection{}\label{exII:SL2} Verify the descriptions of the affine Springer fibers for $G=\SL_{2}$ given in \S\ref{P1 chain} and \S\ref{SL2 ram dim 1}.



\subsubsection{} Let $G=\SL_{2}$ and $\gamma=\mat{0}{t^{n}}{1}{0}\in\frg(F)$. Describe $\sY_{\gamma}$.

\subsubsection{} Let $G=\SL_{2}$ and $\gamma=\mat{0}{t^{n}}{1}{0}\in\frg(F)$. Construct a nontrivial $\Gm$-action on $\sX_{\gamma}$ involving loop rotations (i.e., the action scales $t$) and determine its fixed points.

\subsubsection{} Let $G=\SL_{2}$ and $\gamma=\mat{t^{n}}{0}{0}{-t^{n}}\in\frg(F)$. Let $T\subset G$ be the diagonal torus, then $G_{\gamma}=T\otimes_{k}F$. What is the regular locus $\sX^{\reg}_{\gamma}$ (see \S\ref{dim formula discussion})? Study the $L^+T$-orbits on $\sX_{\gamma}$. 

\subsubsection{} In the setup of \S\ref{sss:L gamma}, show that the action of $\L_{\gamma}$ on $\Gr_{G}$ is free, which implies that its action on $\sX_{\gamma}$ is free. Show also that the permutation action of $\L_{\gamma}$ on the set of irreducible components of $\sX_{\gamma}$ is free.


\subsubsection{} For $G=\GL_{n}$, let $L\subset G$ be the Levi subgroup consisting of block diagonal matrices with sizes of blocks $n_{1},\cdots, n_{s}$, $\sum_{i}n_{i}=n$. Let $\gamma=(\gamma_{1},\cdots,\gamma_{s})\in\frl(F)$ be regular semisimple as an element in $\frg(F)$. What is the invariant $\Delta^{G}_{L}(\gamma)$ (see \eqref{Delta G L}) in terms of familiar invariants of the characteristic polynomials of the $\gamma_{i}$?


\subsubsection{}\label{exII:SL3} Let $G=\SL_{3}$ and $\gamma=\diag(x_{1}t,x_{2}t,x_{3}t)\in\frg(F)$, with $x_{i}\in k$ pairwise distinct and $x_{1}+x_{2}+x_{3}=0$. Describe the affine Springer fibers $\sX_{\gamma}$ and $\sY_{\gamma}$.

Note: this is a good exercise if you have a whole day to kill.

\subsubsection{} For $G=\Sp_{6}$ and $\gamma=\mat{0}{tX}{Y}{0}$ as in Example \ref{ex:Sp}, describe the $\Gm$-fixed points on $\sX_{\gamma}$ and $\sY_{\gamma}$. 



\subsubsection{} Verify the calculations in \S\ref{sss:tilW SL2 ram}.


\subsubsection{} Let $G=\SL_{2}$ and let $\gamma=\mat{t}{0}{0}{-t}$. Describe the affine Springer fiber $\sY_{\gamma}$. What is the action of $\tilW=\jiao{s_{0},s_{1}}$ on $\homog{2}{\sY_{\gamma}}$ in terms of the basis given by the irreducible components of $\sY_{\gamma}$? 



%%%%%%%%%%%%%%%%%%%%%%%%%%%%%


